\section{Homogeneous entries and a cheaper multiplication path}
\label{sec:homogeneous}

The general dense algorithm works for arbitrary inputs. Genuine scanner
entries have more structure, even though the original implementation does
not retain the quantum grading in its object records.

\Needspace{18\baselineskip}
\begin{samepage}
\begin{proposition}[Recovering the erased grading]
\label{prop:grading}
Assign the initial empty object a shift $q=0$. If a parent object of shift
$q$ is extended by smoothing $i\in\{0,1\}$, and that smoothing creates
$c$ closed circles with delooping label $\ell$, assign its child
\begin{equation}
 q'=q+i+c-2\operatorname{popcount}(\ell).
 \label{eq:qshadow}
\end{equation}
For a nonzero differential entry from object $u$ to object $v$, every
supported monomial with $d$ dots satisfies
\begin{equation}
 k(a_u,a_v)-m-2d+q_v-q_u=0,
 \label{eq:qdegree}
\end{equation}
where $2m$ is the current frontier. The invariant holds before and after
each scalar-unit cancellation. In particular, every nonzero entry is
homogeneous in dot degree, and every cancellable entry between equal
matchings is exactly the identity coefficient $1$.
\end{proposition}
\end{samepage}
\begin{proof}
A canonical cobordism with $k$ discs and $d$ dots has degree
$k-m-2d$, obtained by subtracting the arc normalization $m$ from its
Euler degree. A crossing saddle has degree $-1$, compensated by the
smoothing shift $i$. A delooping label on $c$ circles has degree
$c-2\operatorname{popcount}(\ell)$. Its inclusion contributes this
degree and its projection the negative degree. Consequently the source
and target delooping contributions and the smoothing shift cancel exactly
as in~\eqref{eq:qshadow}. This proves~\eqref{eq:qdegree} inductively
through a crossing extension and delooping.

For completeness, composition respects the claimed degree. In a
genus-zero plan without extra dots, Euler characteristic gives
$2r=k_L+k_R+k_O-m$. Every nonzero output term has
$d_O=k_O-r+d_L+d_R$ dots; a closed component obeys the same formula
because its required input dot contributes $-1+1=0$. Substitution shows
\[
 k_O-m-2d_O=(k_L-m-2d_L)+(k_R-m-2d_R).
\]
Thus map degrees add. A homogeneous entry with equal source and target
matching has monomial degrees $-2d$. If its constant term is nonzero,
homogeneity forces $d=0$ for every term, and there is just one such
monomial. It is therefore $1$, and the two object shifts are equal.
The Schur-complement correction in a cancellation is a composition of
homogeneous maps through this identity. Its source and target shifts
telescope, so it has the same degree as the entry to which it is added.
This proves preservation by cancellation.
\end{proof}

This is a recovered grading of the actual ungraded computation. It does
not rely on guessing a degree from the integer encoding. The grading
audit checks every entry of its pre- and post-elimination stage complexes;
transient Schur-complement entries inside a cancellation sequence are
covered by the proof. The implementation still supports arbitrary ungraded
morphisms, so it retains the general ranked-convolution fallback.

\begin{theorem}[Homogeneous convolution]
\label{thm:homogeneous}
Suppose $F,G\in R_r$ are homogeneous of dot degrees $a,b$. Compute their
ordinary subset zeta transforms, multiply the resulting scalar arrays
pointwise, apply inverse zeta transform, and discard every coefficient
whose subset size is not $a+b$. The result is $FG$ in $R_r$. The
transform uses $O((r+1)2^r)$ operations in $\F$.
\end{theorem}
\begin{proof}
Unranked zeta multiplication followed by inversion sums over pairs with
union $S$. Every supported pair has $|A|+|B|=a+b$. Keeping only
$|S|=a+b$ therefore keeps exactly the disjoint pairs. Overlapping pairs
have strictly smaller union and are discarded. The zeta transforms each
perform $r2^{r-1}$ XORs, and the pointwise product and filtering are
linear in $2^r$.
\end{proof}

Contraction preserves the degree of every nonzero monomial. Hence the
theorem applies to contracted genuine scanner entries. The software tests
homogeneity explicitly, so an incorrectly supplied generic input cannot
silently enter this fast path. Neither runtime tests nor
Proposition~\ref{prop:grading} justify applying unranked convolution
without the final cardinality restriction.

\section{Engineering the bit bound}
\label{sec:engineering}

\subsection{Packed integers are not unit-cost coefficients}

There are two distinct uses of integers in the implementation. A subset
mask has at most $m$ bits. A packed morphism can have $2^m$ bits. Repeatedly
clearing its lowest set bit or XORing a shifted output into it can scan a
large fraction of the latter integer on each iteration. Counting those
as constant-time operations would hide an additional exponential factor.

The dense implementation converts a whole morphism to a little-endian
byte array once, reads coefficients by byte and bit offset, accumulates
output coefficients in a byte array, and converts the result back once.
Rank polynomials in the general convolution are integers of only
$r+1$ bits. Carryless multiplication uses shifts and XORs on these small
polynomials; ordinary integer multiplication would introduce carries and
would be mathematically wrong.

\subsection{Adaptive sparse and dense composition}

Dense transforms have setup costs, and most entries in small diagrams
are sparse. The adapter therefore chooses between an explicit sparse
evaluation and the dense factorization. Correctness does not depend on
which branch is selected.

Let $s_F,s_G$ be the support sizes. If component $s$ has $b_s$ output
circles, a supported monomial pair expands to at most
\[
 D=\prod_s\max(1,b_s)
\]
output monomials. The sparse cost estimate includes
\[
 s_Fs_G\bigl(r+1+(k_O+1)D\bigr),
\]
not just $s_Fs_G$. It accounts for output expansion and for the length of
small monomial masks. The implementation compares this estimate to a
fixed multiple of
\[
 (k_L+1)2^{k_L}+(k_R+1)2^{k_R}
 +(k_O+1)(r+1)2^{k_O}+(r+1)^2 2^r.
\]
The sparse branch toggles individual coefficient bytes instead of
repeatedly XORing full morphism integers. Its selected work therefore
differs from the dense envelope by at most a polynomial factor in the
frontier size. The multiplier is a performance parameter held fixed in
the asymptotic assertion.

This design also removes the baseline's potentially enormous dictionary
of individual monomial-pair evaluations for a dense composition. It caches
the factorization plan, whose storage is proportional to coefficient-table
size up to a polynomial factor. Retaining many distinct plans can still
consume substantial memory; a per-plan bound is not a bound on the number
of plans.

\subsection{Crossing transfers}

Faster composition alone would leave another packed-integer loop in
crossing transfer. The transfer compiler already isolates the input
circles touched by the four crossing slots. There are at most four such
input circles, together with at most two local smoothing discs. The local
surface has a universally bounded number of output circles and at most
sixteen source/target delooping-label pairs. For example, if $V\leq6$ is
the number of local discs and $E\leq4$ the number of local gluings, then
the sum of local genus-zero boundary counts is at most $V+E\leq10$.
Capping closed components can only reduce that boundary count.

The byte-transfer path memoizes the local output monomial indices for
each touched input pattern. It renumbers untouched dots and toggles the
corresponding output coefficient bytes. The local expansion factor is a
constant, while renumbering costs a polynomial in the frontier size. A
fixed small-support fallback retains the old routine; with its support
bounded by a fixed constant, even full-integer updates fit within
$\poly(w)2^{w/2}$. Forced-mode tests exercise the byte path even on those
small cases.

Geometry setup must be charged separately. Even when two overlay
matchings have only one circle, walking their arcs can require time
proportional to the frontier. Thus a coefficient bound expressed in
$k_L,k_R,k_O,r$ presupposes an already compiled plan. The full-call bound
uses the frontier $w$ and the bit length of its labels.

\section{What the local bound implies for the whole scanner}
\label{sec:global}

Let $N$ be the largest number of physically stored objects at any stage,
including objects about to be cancelled, and let $w$ be the largest exposed
frontier. These are execution parameters; neither is assumed to be small.
Repeated-summand weights, if another backend uses them, are additional
data whose arithmetic must also be charged.

\begin{theorem}[Scanner resource envelope]
\label{thm:execution}
For the finite characteristic-two scanner with the coefficient routines
above, standard explicit object and differential storage, and ordinary
unit-cancellation bookkeeping, the running time is bounded by
\begin{equation}
 \poly(n,N,w,\log(n+2))\,2^{w/2}.
 \label{eq:envelope}
\end{equation}
The bound allows polynomial factors for dictionary lookup or a deterministic
comparison-map replacement. Equivalently, if $Q$ counts all executed
coefficient operations and primitive data-structure probes, key comparisons,
and queue moves, a more informative
execution bound is
\begin{equation}
 \poly(n,w,\log(N+2))\,(Q+1)2^{w/2}.
 \label{eq:executed}
\end{equation}
No inference from boundary width alone to $N$ or $Q$ is made.
\end{theorem}
\begin{proof}
A stage has at most $N^2$ possible differential entries. One crossing
has only a constant number of smoothings and delooping choices per object,
so allocating objects and transferring entries require polynomially many
operations in $N$. A unit cancellation removes two objects and updates at
most $N^2$ entries. There are at most $N/2$ such cancellations per stage.

For a simple implementation that rescans all possible pivots, this already
gives polynomial control. The shipped lazy min-fill queues also have
polynomial control. There are polynomially many initial or updated offers;
a stale offer is reinserted only with a strictly larger integer cost, and
the cost lies between $0$ and $(N-1)^2$. Thus even the deliberately loose
bound obtained by charging every offer all possible promotions is
polynomial in $N$. Plan caches contain no more entries than operations
that create them; the interning tables across all $n$ stages also have
polynomially many records in $n,N$. Collision resolution or balanced-tree
lookup changes only that polynomial bound.

Every composition and transfer has cost $\poly(w,\log(n+2))2^{w/2}$.
An addition or serialization of a morphism has no larger cost. Object
identifiers and degree indices have polynomially bounded bit lengths.
Multiplying the operation count by these costs proves both formulas.
This proof does not attempt to optimize the polynomial exponent of the
bookkeeping.
\end{proof}

\Needspace{12\baselineskip}
\begin{samepage}
\begin{corollary}[A sufficient execution regime]
\label{cor:qp}
If a family of executions satisfies
\[
 w+\log(N+2)=O((\log(n+2))^2),
\]
then the improved scanner runs in $n^{O(\log n)}$ time on that family.
The same conclusion follows from
$w+\log(Q+2)=O((\log(n+2))^2)$ in~\eqref{eq:executed}.
\end{corollary}
\end{samepage}
\begin{proof}
Take the logarithm of the right-hand side of the relevant resource
envelope. Each polynomial factor contributes
$O(\log(n+2)+\log(N+2)+\log(w+2))$ or the corresponding term in $Q$;
the coefficient-table factor contributes $O(w)$. The sum is
$O((\log(n+2))^2)$.
\end{proof}

The corollary gives a precise target for a structural theorem. It is not
such a structural theorem. A small frontier bounds the coefficient
dimension of one morphism, but does not bound how often the same matching
must occur. If $\mu$ denotes a bound on multiplicity per matching and
degree, a safe unrestricted matching count on $w$ labeled points is
$(w-1)!!$. Consequently one obtains, at best from that information alone,
\[
 \log N=O\bigl(\log(n+2)+w\log(w+1)+\log(\mu+1)\bigr).
\]
A Catalan bound requires an actual certified disk order with noncrossing
matchings. The prior continuation supplies counterexamples to using such
a bound for arbitrary disconnected frontiers, so we do not reuse it here.

\subsection{A sharper bound for component-sharing types}
\label{sec:types}

The recovered grading also improves the finite-type estimate used by the
earlier structural-compression report. This is an improvement to a
sufficient parameter regime, beyond the faster composition kernel.

For integers $m\geq0$ and $b\geq1$, let
\[
 B_m=\binom{m}{\floor{m/2}},\qquad M(2m)=(2m-1)!!,
\]
with $M(0)=1$. A homogeneous element of a $k$-circle coefficient space
belongs to one of its $k+1$ dot-degree layers. Each layer has dimension
at most $B_m$ for $k\leq m$. Hence the number of possible coefficients,
including zero, is at most $H_m=(m+1)2^{B_m}$. This union bound
intentionally overcounts zero.

\begin{theorem}[Homogeneous component-type bound]
\label{thm:types}
At a fixed frontier of $2m$ labeled points, the number of serialized
connected differential components with at most $b$ objects, homogeneous
entries, and homological degrees normalized by their minimum is at most
\begin{equation}
 K(m,b)=\sum_{s=1}^{b}s^s M(2m)^s H_m^{s^2}.
 \label{eq:types}
\end{equation}
Consequently
\begin{equation}
 \log K(m,b)=O\bigl(b^2 B_m+b m\log(m+1)+b\log(b+1)\bigr)
 =O(b^2 B_m).
 \label{eq:typecost}
\end{equation}
\end{theorem}
\begin{proof}
For a connected component of $s$ objects, every nonzero differential
changes homological degree by one. An undirected path of at most $s-1$
edges joins a minimum-degree vertex to any other vertex. The normalized
degrees therefore lie in $\{0,\ldots,s-1\}$, giving at most $s^s$
assignments. There are at most $M(2m)^s$ choices of matching labels. For
each ordered pair of vertices, the homogeneous coefficient belongs to a
set of at most $H_m$ possibilities. This gives $H_m^{s^2}$ choices.
We have ignored the required degree difference, connectedness, and
$d^2=0$, so the count is an upper bound. It already counts ordered
serializations; no graph-isomorphism algorithm is assumed.

Take logarithms and bound the sum by $b$ times its largest crude term.
The estimates $\log M(2m)=O(m\log(m+1))$ and the definition of $H_m$
give the first expression. Finally, $m\log(m+1)=O(B_m)$,
$\log(m+1)=O(B_m)$, and $\log(b+1)=O(b)$ uniformly over nonnegative
integers, with the finitely many small values absorbed into constants.
This gives the second expression.
\end{proof}

The keys in the earlier component-sharing code do not store quantum
shifts. This does not invalidate the count: Proposition~\ref{prop:grading}
ensures that each observed coefficient is homogeneous. Forgetting a
hidden grading cannot increase the number of structures compared with
the set of possible graded structures. There is also a potentially
tighter count that assigns normalized quantum shifts to vertices.
Across a nonzero edge their difference is
$m-k+2d\in[0,2m]$. In a connected $s$-vertex component their range is
therefore at most $2m(s-1)$. Once both matching and quantum shifts are
fixed, every ordered pair has a uniquely determined dot-degree layer.
This yields the alternative upper bound
\[
 \sum_{s=1}^{b}s^s[2m(s-1)+1]^s M(2m)^s 2^{s^2 B_m}.
\]
Its leading exponent has the same order as~\eqref{eq:typecost}.

\begin{corollary}[Improved sufficient regime for exact sharing]
\label{cor:sharing}
Consider the exact component-sharing scanner from report 07, using the
present coefficient kernel. Suppose every reduced connected component
encountered in an execution has at most $b$ objects and the frontier has
at most $2m$ points. Then a sufficient condition for
$n^{O(\log n)}$ time is
\begin{equation}
 b^2\binom{m}{\floor{m/2}}=O((\log(n+2))^2).
 \label{eq:sharingregime}
\end{equation}
\end{corollary}
\begin{proof}
There are at most $K(m,b)$ stored component representatives, each with
at most $b$ objects. Tensoring a representative with one crossing creates
only a constant factor more objects before cancellation. Distinct parent
summands stay independent under this operation, so the representation
has polynomial size in $bK(m,b)$. Its homological-shift weights have
degree $O(n)$ and coefficients of $O(n)$ bits: the full cube has at most
exponentially many generators in $n$. Saturated decision weights are
smaller still. The resource envelope is therefore polynomial in
$n,b,K(m,b)$ times $\poly(m)2^m$. Insert
Theorem~\ref{thm:types} and~\eqref{eq:sharingregime}.
\end{proof}

Since $B_m=\Theta(2^m/\sqrt{m+1})$, the sufficient leading term improves
from $b^2 2^m$ in the generic ungraded type count to
$b^2 2^m/\sqrt{m+1}$. For bounded $b$, put
$L=\log_2(n+2)$ and $\ell=\log_2 L$. The corresponding frontier scale
is enlarged from $w\leq4\ell+O(1)$ to
\[
 w\leq4\ell+\log_2\ell+O(1).
\]
The additive constant can depend on the chosen bound for $b$.
This is a modest asymptotic refinement of the earlier regime. The
unresolved task is still to prove useful bounds on actual connected
component sizes for broad knot families. The archive does not claim a
new benchmark of the report-07 scanner with every recent branch merged.


\subsection{Extremal slices and geometric alternatives}

Kelom\"aki and Sch\"utz give an independent useful route for extremal
homological slices~\cite{kelomaki}. Their prescribed nice scans and Morse
cancellations control formal-object multiplicities by binomial coefficients.
Tracking their parameters yields a slice envelope of the form
$2^{O(g)}(g+1)^{O(1)}n^{O(k+1)}$, with endpoint girth $g$ and extremal
depth $k$ from a chosen endpoint of the unshifted homological cube
grading. Raw shallow-cube enumeration does not justify this bound: it
can retain exponentially many circle labels. The specified cancellations
are necessary, and a nice scan of small girth must be supplied.

The dense coefficient kernel could be used in the $\F$ specialization of
that construction, but
this archive does not implement or benchmark that integration. Nor do
extremal slices automatically detect every nontrivial knot: the necessary
homology may occur outside the selected slice. On the geometric route,
the analogous missing work is constructive control of hierarchy search,
compressed cutting, and primitive bit lengths. Both routes are worth
pursuing, with their construction and verification costs stated separately.
